\documentclass[11pt]{article}

\usepackage[margin=1in]{geometry}
\usepackage{hyperref}

\setlength{\parindent}{0pt}

\title{Station-Level Bike Demand Prediction in Toronto: \\ 
A Machine Learning Model Comparison Study}

\author{}
\date{}

\begin{document}

\maketitle


\section{Introduction}

Bike-sharing systems are an important part of urban transportation in many cities.
In Toronto, the Bike Share Toronto program provides bicycles across hundreds of
stations and supports a large number of daily trips. However, bike demand varies
significantly depending on time and location. Predicting bike demand can help
operators improve bike allocation and reduce shortages at busy stations.

In this project, we aim to predict \textbf{bike demand at specific stations and
    times} using historical ridership data from Toronto Bike Share. The task is
formulated as a \textbf{regression problem}, where the goal is to predict the
number of bike trips starting from a station during a given hour.

Our research question is:

\begin{quote}
    How do different machine learning models compare in predicting station-level bike demand?
\end{quote}


\section{Dataset}

We use the \textbf{Toronto Bike Share Ridership Data (2021–2024)} available on Kaggle:

\begin{center}
    \url{https://www.kaggle.com/datasets/mattskogstadstubbs/toronto-bike-share-ridership-data-2021-2024}
\end{center}

The dataset contains trip-level records including start time, start station, trip
duration, bike ID, and user type, covering multiple years of ridership data. We convert the trip-level data into station-level
demand data by aggregating trips by \textbf{station and hour}. The target variable
is defined as the number of trips starting from a station within a given hour.

\section{Feature Engineering}

Several features will be constructed to help the models capture temporal and spatial
demand patterns.

Temporal features extracted from trip start time include hour of day, day of week,
month, and whether the day is a weekend. These features capture daily and weekly
usage cycles.

Station-level features include the station ID and historical average demand for
each station. In addition, we will compute the proportion of annual members among
users as a behavioral feature reflecting commuter usage patterns.

We will also include lag features such as the demand in the previous hour and the
demand at the same hour on the previous day, which help capture temporal dependence
in bike demand.

\section{Models}

We will conduct a comparison study using several machine learning models.

Classical models include:

\begin{itemize}
    \item Linear Regression
    \item Ridge Regression
    \item Random Forest Regression
\end{itemize}

As an advanced model not covered in the lecture, we will implement:

\begin{itemize}
    \item XGBoost (Extreme Gradient Boosting)
\end{itemize}

\section{Evaluation}

The models will be evaluated using standard regression metrics such as Mean Absolute
Error (MAE), Mean Squared Error (MSE), and Root Mean Squared Error (RMSE). The
dataset will be split using a time-based train-test split to avoid information leakage.

\section{Originality}

As an extension, we will cluster bike stations based on their historical demand
patterns using K-means clustering. The resulting station cluster labels will be
used as additional features to capture usage patterns among different types of
stations.

\end{document}